\chapter{Dynamik}

Viele aus dem Maschinenbau bekannte Bauteile, wie etwa Kurbelwelle, Kolben oder Pleuel
bewegen sich beschleunigt und sind dynamisch beansprucht. Untersucht man diese
Bewegungen im Zusammenhang mit den vorhandenen Kräften, so spricht man von Dynamik.

\begin{figure}[H]    
    \centering          %     left botm right top
    \includegraphics[clip, trim=0cm 0cm 0cm 0cm,width=.4\textwidth]{einteilungMechanik.png} %.6 stands for 0.6 x textwidth
    \caption{Einteilung Mechanik \cite{bartsch_dynamik_2023}}
    \label{fig:einteilungMechanik_other}
\end{figure}

In der Statik und Festigkeitslehre werden im Wesentlichen Systeme betrachtet, die in Ruhe
sind; die Zeit als eine Grundgrösse der Mechanik kommt also nicht vor. In der Dynamik spielt
die Zeit hingegen eine entscheidende Rolle.

Die 3 Grundgesetze von Newton lauten:

\begin{enumerate}
    \item Jeder Körper verharrt in seinem Zustand der gleichförmigen, geradlinigen Bewegung, 
    wenn er nicht durch einwirkende Kräfte gezwungen wird, diesen Zustand zu ändern
    \item Die Änderung der Bewegungsgrösse ist proportional zur einwirkenden Kraft
    \item Die Wirkung zweier Körper aufeinander ist stets gleich und von entgegengesetzter 
    Richtung (actio = reactio)
\end{enumerate}

Insbesondere das 2. Gesetz ist für die Dynamik zentral; mathematisch formuliert (rechts mit Produktregel):

\begin{minipage}{.45\textwidth}
    \begin{equation}
        \overrightarrow{F} = \frac{d(m \cdot \overrightarrow{v})}{dt}
    \end{equation}
\end{minipage}
\hspace{.1\textwidth}
\begin{minipage}{.45\textwidth}
    \begin{equation}
        \overrightarrow{F} = m \cdot \frac{d \overrightarrow{v}}{dt} + \frac{dm}{dt} \cdot \overrightarrow{v}
    \end{equation}
\end{minipage}

Für den starren Körper ist der zweite Term gleich Null und es bleibt die unter dem Namen
Newton'sches Gesetz bekannte Formel \cite{bartsch_dynamik_2023}:

\begin{equbox}
    \begin{equation}
        \overrightarrow{F} = m \cdot \overrightarrow{a}
    \end{equation}
\end{equbox}

Auf Basis der obigen Grundgesetze sind folgende Begriffe und Formeln aus der Physik
bekannt. Analog zu einer Translation mit der Grösse Masse m existieren die entsprechenden
Beziehungen für eine Rotation um eine feste Achse mit der Grösse Massenträgheitsmoment $J$.

\begin{table}[h]
    \centering
    \caption{Analogie dynamische Grössen Translation und Rotation \cite{bartsch_dynamik_2023}}
    \label{tab:anal_dyn_grössen}
    \begin{figure}[H]    
        \centering          %     left botm right top
        \includegraphics[width=\textwidth]{analogie_dynamische_grössen.png} %.6 stands for 0.6 x textwidth
    \end{figure}
\end{table}

\section{Newton oder D'alembert}
\label{sec:newtondalembert}

Mit Hilfe des Grundgesetzes können entweder bei gegebenen Kräften die Bewegung eines
Bauteils oder bei gegebener Bewegung die auftretenden Kräfte bestimmt werden. Dabei sind
nach Newton mit $\overrightarrow{F}$ die Summe aller äusseren Kräfte, also die resultierende Kraft gemeint \cite{bartsch_dynamik_2023}.

\begin{bspbox}
    \newline
    \begin{minipage}{.55\textwidth}
        
        \textbf{Newton:} Ein Paket der Masse $m = 50 kg$ gleitet auf einer Rutsche
        nach unten. Die Gleitreibungszahl beträgt $\mu = 0,4$.
        Wie gross ist die Endgeschwindigkeit bei einer
        Anfangsgeschwindigkeit von $v_0 = 2 ms^{-1}$? \newline

        Kräftegleichgewicht in y-Richtung:

        \begin{align*}
            \sum F_y &= N - m \cdot g \cdot \cos(30^\circ) = 0 \\
            \rightarrow N &= m \cdot g \cdot \cos(30^\circ)
        \end{align*}

        Die resultierende Kraft in x-Richtung beträgt:

        \begin{equation*}
            F_{res,x} = m \cdot g \cdot \sin(30^\circ) - R
        \end{equation*}

    \end{minipage}
    \hspace{.05\textwidth}
    \begin{minipage}{.38\textwidth}
        \begin{figure}[H]
            \centering
            \begin{subfigure}{\textwidth}
                \includegraphics[clip, trim=0cm 0cm 0cm 0cm, width=\textwidth]{newtonbsp_1.png}
                \label{fig:newtonbsp1}
            \end{subfigure}
        \end{figure}
    \end{minipage}

    \begin{minipage}{.55\textwidth}

        Mit dem Reibungsgesetz $R = \mu \cdot N$ erhalten wir:

        \begin{align*}
            F_{res, x} &= m \cdot g \cdot (\sin(30^\circ) - \mu \cdot \cos(30^\circ)) = m \cdot a \\
            \rightarrow a &= g \cdot (\sin(30^\circ) - 0.4 \cdot \cos(30^\circ)) = 1.51ms^{-2}
        \end{align*}

        Mit der kinematischen Beziehung:

        \begin{equation*}
            v^2 = 2 \cdot a \cdot (x - x_0) + v_0^2
        \end{equation*}

        Ergibt sich die Endgeschwindigkeit:

        \begin{align*}
            v &= \sqrt{2 \cdot a \cdot (x - x_0) + v_0^2} \\
            &= \sqrt{2 \cdot 1.51ms^{-2} \cdot 2m + (2ms^{-1})^2} = 3.17ms^{-1}
        \end{align*}

    \end{minipage}
    \hspace{.05\textwidth}
    \begin{minipage}{.38\textwidth}
        \begin{figure}[H]
            \centering
            \begin{subfigure}{\textwidth}
                \includegraphics[clip, trim=0cm 0cm 0cm 0cm, width=\textwidth]{newtonbsp_2.png}
                \label{fig:newtonbsp2}
            \end{subfigure}
        \end{figure}
    \end{minipage}
    \newline
\end{bspbox}

Im Gegensatz zu Newton fasst D'alembert den Ausdruck $m \cdot \overrightarrow{a}$ als Trägheitskraft auf.
Damit lässt sich ein dynamisches Problem auf ein statisches Gleichgewicht reduzieren \cite{bartsch_dynamik_2023}.

\begin{bspbox}
    \begin{minipage}{.55\textwidth}

        \textbf{D'alembert:} In der Skizze des freigemachten Bauteils wird die
        Trägheitskraft $m \cdot a$ entgegengesetzt der Bewegung
        eingeführt. \newline

        Das Gleichgewicht für beide Richtungen lautet:

        \begin{align*}
            \sum F_y &= N - m \cdot g \cdot \cos(30^\circ) = 0 \\
            \sum F_x &= m \cdot g \cdot \sin(30^\circ) - R - m \cdot a = 0
        \end{align*}

    \end{minipage}
    \hspace{.05\textwidth}
    \begin{minipage}{.38\textwidth}
        \begin{figure}[H]
            \centering
            \begin{subfigure}{\textwidth}
                \includegraphics[clip, trim=0cm 0cm 0cm 0cm, width=\textwidth]{bspdalembert.png}
                \label{fig:dalembertbsp}
            \end{subfigure}
        \end{figure}
    \end{minipage}
\end{bspbox}

Mit dem Vorgehen nach D'alembert führen wir ein dynamisches Problem formal auf ein
statisches Problem zurück. Dies mag ein Grund dafür sein, dass es in der Praxis oft bevorzugt
angewendet wird.

\section{Energieerhaltung}

Mit dem aus der Physik bekannten Begriff der Arbeit kann man auf Basis
des dynamischen Grundgesetzes die kinetische Energie ableiten. Neben der kinetischen
Energie existieren noch die potentielle- und die elastische Energieformen in der Mechanik \cite{bartsch_dynamik_2023}.

Die kinetische Energie eines translatorisch bewegten starren 
Körpers lässt sich schreiben (rechts für rotierende Körper, Massenträgheitsmomente
für Grundkörper sind im Anhang \ref{app:C_tabellen}, Tabelle \ref{tab:Massenträgheitsmomente} 
zu finden):

\begin{equbox}
    \begin{minipage}{.44\textwidth}
        \begin{equation}
            E_{kin} = \frac{1}{2} \cdot m \cdot v^2
        \end{equation}
    \end{minipage}
    \hspace{.1\textwidth}
    \begin{minipage}{.44\textwidth}
        \begin{equation}
            E_{kin} = \frac{1}{2} \cdot J \cdot \omega^2
        \end{equation}
    \end{minipage}
\end{equbox}

Die potentielle Energie bezüglich eines festgelegten Niveaus berechnet sich wie folgt:

\begin{equbox}
    \begin{equation}
        E_{pot} = m \cdot g \cdot h
    \end{equation}
\end{equbox}

Federn können elastische Energie speichern, sie beträgt für lineare Kennlinien:

\begin{minipage}{.54\textwidth}
    \begin{equbox}
        \begin{equation}
            E_{feder} = \frac{1}{2} \cdot c \cdot (x_2^2 - x_1^2)
        \end{equation}
    \end{equbox}
    \vspace{10pt}
    Elastische Bauteile können ebenfalls als Federn
    aufgefasst werden und sind analog zu behandeln. So
    können Zugstäbe, Biegebalken und Torsionsstäbe
    elastische Energie speichern \cite{bartsch_dynamik_2023}.
\end{minipage}
\hspace{.1\textwidth}
\begin{minipage}{.34\textwidth}
    \begin{figure}[H]   
        \centering          %     left botm right top
        \includegraphics[clip, trim=0cm 0cm 0cm 0cm,width=\textwidth]{Efeder.png} %.6 stands for 0.6 x textwidth
    \end{figure}
\end{minipage}

Bei Bewegungsvorgängen können diese drei Energieformen ineinander überführt werden.
Unter Berücksichtigung der Verluste bleibt dabei die Summe der Energien für zwei 
Zustände (1 und 2) konstant:

\begin{equbox}
    \begin{equation}
        \sum E_1 - W_R = \sum E_2
    \end{equation}
\end{equbox}

Mit $\sum E_i$ sind die drei Energieformen kinetische, potentielle und elastische Energie gemeint. $W_R$ sind die 
Verluste. Wird einem System Energie zugeführt, z.B.
mittels Antrieb, so ist die zugeführte Arbeit auf der linken Seite zu addieren \cite{bartsch_dynamik_2023}.

\begin{bspbox}
    \begin{minipage}{.55\textwidth}
        Das Beispiel aus Kapitel \ref{sec:newtondalembert} wird alternativ mit
        Hilfe des Energieerhaltungssatzes gelöst.
        Zunächst wird ein Referenzniveau für die
        potentielle Energie festgelegt (Epot = 0).
        Die Verluste infolge Reibung betragen bei
        konstanter Reibkraft über dem Weg ∆𝑥 = 2m:

        \begin{equation*}
            W_R = R \cdot \Delta x = \mu \cdot m \cdot g \cdot \cos(30^\circ) \cdot \Delta x
        \end{equation*}

    \end{minipage}
    \hspace{.05\textwidth}
    \begin{minipage}{.38\textwidth}
        \begin{figure}[H]
            \centering
            \begin{subfigure}{\textwidth}
                \includegraphics[clip, trim=0cm 0cm 0cm 0cm, width=\textwidth]{energieerhaltungbsp.png}
                \label{fig:newtonbsp1}
            \end{subfigure}
        \end{figure}
    \end{minipage}

    Der Energieerhaltungssatz als Bilanz zwischen zwei Zuständen „oben“ und „unten“ lautet:

    \begin{equation*}
        \frac{1}{2} \cdot m \cdot v_0^2 + m \cdot g \cdot h - W_R 
        = \frac{1}{2} \cdot m \cdot v^2 \rightarrow v 
        = \sqrt{v_0^2 + 2 \cdot g \cdot (h - \mu \cdot \cos(30^\circ) \cdot \Delta x)} = 3.17ms^{-1}
    \end{equation*}

\end{bspbox}

\section{Impuls- und Drehimpulserhaltung}

Die von Newton eingeführte Bewegungsgrösse $m \cdot \overrightarrow{v}$ wird Impuls $\overrightarrow{p}$ 
genannt. 

\subsection{Impulssatz}

Wir gehen aus vom dynamischen Grundgesetz für konstante Massen und 
integrieren über die Zeit. Damit erhalten wir den Impulssatz in integraler Form:

\begin{equbox}
    \begin{equation}
        m \cdot \overrightarrow{v_1} + \int \overrightarrow{F} dt = m \cdot \overrightarrow{v_2}
    \end{equation}
\end{equbox}

Das Integral wird Kraftstoss genannt, die Änderung des Impulses ist gleich dem Kraftstoss.
Greifen mehrere Kräfte an, so ist für $\overrightarrow{F}$ die resultierende Kraft einzusetzen.

Der Impulssatz kann für einen Massenpunkt wie folgt gedeutet werden. Der Einfachheit
halber betrachten wir eine geradlinige Bewegung und lassen deshalb die Vektorpfeile weg.

\begin{figure}[H]    
    \centering          %     left botm right top
    \includegraphics[clip, trim=0cm 0cm 0cm 0cm,width=.7\textwidth]{impulssatz.png} %.6 stands for 0.6 x textwidth
    \caption{Impulssatz \cite{bartsch_dynamik_2023}}
    \label{fig:impulssatz}
\end{figure}

Eine Masse bewegt sich mit der Geschwindigkeit $v_1$ nach rechts. Es greift über eine (kurze)
Zeit eine Kraft $F$ an, die nicht konstant sein muss. Durch diesen Kraftstoss $\int F dt$ verändert
sich die Geschwindigkeit und beträgt nachher $v_2$.

Die Stärke des Impulssatzes liegt in der Anwendung für schnelle Vorgänge wie Stösse oder
Kollisionen, bei welchen die Stosszeit Bruchteile von Sekunden beträgt.
Er gilt aber auch für „langsame“ Vorgänge; dies wird anhand des einfachen Beispiels gezeigt.

\begin{bspbox}
    \begin{minipage}{.55\textwidth}
        Das Beispiel aus Kapitel \ref{sec:newtondalembert} wird alternativ mit
        Hilfe des Impulssatzes gelöst.

        In diesem Beispiel ist die resultierende Kraft über der Zeit des
        Kraftstosses konstant und beträgt:

        \begin{equation*}
            F_{res} = m \cdot g \cdot (\sin(30^\circ) - \mu \cdot \cos(30^\circ))
        \end{equation*}

        Das Integral $\int F dt = F_{res} \cdot \Delta t$ und der Impulssatz:

    \end{minipage}
    \hspace{.05\textwidth}
    \begin{minipage}{.38\textwidth}
        \begin{figure}[H]
            \centering
            \begin{subfigure}{\textwidth}
                \includegraphics[clip, trim=0cm 0cm 0cm 0cm, width=\textwidth]{newtonbsp_1.png}
                \label{fig:newtonbsp}
            \end{subfigure}
        \end{figure}
    \end{minipage}

    \begin{align*}
        m \cdot v_0 + F_{res} \cdot \Delta t &= m \cdot v \\
        m \cdot v_0 + m \cdot g \cdot (\sin(30^\circ) - \mu \cdot \cos(30^\circ)) \cdot \Delta t &= m \cdot v \\
        g \cdot (\sin(30^\circ) - \mu \cdot \cos(30^\circ)) = \frac{v - v_0}{\Delta t} = a
    \end{align*}

    Mit der kinematischen Beziehung $v^2 = 2 \cdot a \cdot (s - s_0) + v_0^2$ ergibt sich die Endgeschwindigkeit.
\end{bspbox}

Der Impulssatz lässt sich auch für Systeme aus mehreren Massen oder für starre Körper
anwenden. Der Massenmittelpunkt oder Schwerpunkt bewegt sich so, als wäre die gesamte
Masse in diesem Punkt vereinigt und als würden alle Kräfte in ihm angreifen.

\subsection{Impulserhaltungssatz}

Verschwindet bei einem abgeschlossenen System die resultierende Kraft, so wird der
Impulssatz mit $\overrightarrow{F} = 0$ zum Impulserhaltungssatz.

\begin{equbox}
    \begin{equation}
        \sum (m_i \cdot \overrightarrow{v_i})_1 = \sum (m_i \cdot \overrightarrow{v_i})_2
    \end{equation}
\end{equbox}

Dieser Sachverhalt gilt für ein System von mehreren Massenpunkten oder für einen starren
Körper. Das bedeutet, dass der Massenmittelpunkt von mehreren Massenpunkten oder der
Schwerpunkt eines starren Körpers sich gleichförmig geradlinig bewegt oder in Ruhe bleibt
(Spezialfall), wenn keine äusseren Kräfte angreifen.

Der Impulserhaltungssatz gilt ohne Einschränkungen, d.h. auch bei Vorgängen mit
„Energieverlusten“ wie Reibung oder plastischer Deformation. Stossvorgänge, die in der
Praxis immer mit Verlusten verlaufen, sind deshalb nur über den Impulserhaltungssatz
erfassbar.

\subsection{Drehimpuls}

In Analogie zum Impuls $m \cdot \overrightarrow{v}$ bei Translationen existiert für Rotationen starrer Körper die
Bewegungsgrösse $J \cdot \overrightarrow{\omega}$, genannt Drehimpuls $\overrightarrow{L}$.

\begin{equation}
    \overrightarrow{L} = J \cdot \overrightarrow{\omega}
\end{equation}

Zunächst soll der Spezialfall einer Rotation um eine Symmetrieachse betrachtet werden; dann
sind $\overrightarrow{L}$ und $\overrightarrow{\omega}$ gleichgerichtet. Da der Schwerpunkt auf der Symmetrieachse liegt, wird das
Massenträgheitsmoment bezüglich dem Schwerpunkt $J_S$ verwendet.
In Analogie zum Impulssatz existiert für Rotationen der Drehimpulssatz:

\begin{equation}
    \overrightarrow{M} = J_S \cdot \frac{d \overrightarrow{\omega}}{dt}
\end{equation}

Besser bekannt unter dem Namen \textbf{Drallsatz} in folgender Schreibweise:

\begin{equbox}
    \begin{equation}
        \overrightarrow{M} = J_S \cdot \overrightarrow{\alpha}
    \end{equation}
\end{equbox}

In integraler Form:

\begin{equbox}
    \begin{equation}
        J_S \cdot \overrightarrow{\omega}_1 + \int \overrightarrow{M}_S dt = J_S \cdot \overrightarrow{\omega}_2
    \end{equation}
\end{equbox}

Die Änderung des Drehimpulses ist gleich dem Zeitintegral des Momentes $\int M dt$.
Der Drehimpulssatz kann wie folgt gedeutet werden. Die Vektoren $\overrightarrow{M}$ und $\overrightarrow{\omega}$ sind parallel und
wir lassen deshalb die Vektorpfeile weg.

\begin{figure}[H]    
    \centering          %     left botm right top
    \includegraphics[clip, trim=0cm 0cm 0cm 0cm,width=.7\textwidth]{drehimpulssatz.png} %.6 stands for 0.6 x textwidth
    \caption{Drehimpulssatz \cite{bartsch_dynamik_2023}}
    \label{fig:drehimpulssatz}
\end{figure}

Ein Körper dreht sich mit der Winkelgeschwindigkeit 𝜔1. Es greift über eine (kurze) Zeit ein
Moment $M$ an, das nicht konstant sein muss. Durch den Drall $\int M dt$ verändert sich die
Winkelgeschwindigkeit und beträgt nachher $\omega_2$.

Der Drehimpulssatz oder Drallsatz wird meist in der Form $\overrightarrow{M} = J \cdot \overrightarrow{\alpha}$ verwendet, wobei die
Grösse $J \cdot \overrightarrow{\alpha}$ als „Trägheitskraft“ nach D'alembert eingeführt wird.

Liegt die Drehachse schief zur Symmetrieachse des Bauteils, so zeigt der Vektor des
Momentes $\overrightarrow{M}$ nicht mehr in Richtung der Drehachse $\overrightarrow{\omega}$. Dies zeigt sich beispielsweise bei
dynamisch nicht ausgewuchteten Rotoren

\subsection{Drehimpulserhaltungssatz}

Der Drehimpulserhaltungssatz ist wie der Impulserhaltungssatz universell gültig, also auch
bei Vorgängen mit Reibung und plastischer Deformation, bei denen mechanische Energien in
Wärme umgewandelt werden. Er ist deshalb zur Erfassung von Stössen unabdingbar.

\begin{equbox}
    \begin{equation}
        \sum (J_i \cdot \overrightarrow{\omega}_i)_1 = \sum (J_i \cdot \overrightarrow{\omega}_i)_2
    \end{equation}
\end{equbox}

Eine andere Anwendung stellt das Kuppeln von rotierenden Wellen dar. Liegen bei einem
Getriebe mehrere Wellen parallel, so wirken beim Kuppeln Umfangskräfte und Lagerkräfte,
die als äusseres Moment aufgefasst werden können. Somit gilt der Drehimpulserhaltungssatz
nur bei fluchtenden Wellen. Man kann allerdings das Problem bei Getrieben umgehen, indem
man die Massenträgheitsmomente auf die zwei fluchtenden Wellen reduziert.

\begin{bspbox}
    Ein Getriebe besteht aus vier Wellen, mit den Übersetzungen $i_{AB} = 3$ und $i_{CD} = 4$. Zwischen den Wellen
    B und C ist eine Kupplung, die zunächst offen ist. Die Wellen B und C haben den gleichen Drehsinn,
    aber unterschiedliche Drehzahlen. Wie gross ist die Drehzahl der Welle BC nach dem Kuppeln und
    wie gross ist die in der Kupplung entstehende Wärme?

    Die auf die Achse B bzw. C reduzierten Trägheitsmomente:

    \begin{minipage}{.55\textwidth}
        \begin{align*}
            J_{redB} &= J_B + J_A \cdot (\frac{n_A}{n_B})^2 \\
            &= 10kgm^2 + 270kgm^2 \cdot (\frac{1s^{-1}}{3s^{-1}}) = 40kgm^2 \\
            J_{redC} &= J_C + J_D \cdot (\frac{n_D}{n_C})^2 \\
            &= 15kgm^2 + 1kgm^2 \cdot 4 = 31kgm^2
        \end{align*}

        Für das reduzierte System gilt der Drehimpulserhaltungssatz:

    \end{minipage}
    \hspace{.05\textwidth}
    \begin{minipage}{.38\textwidth}
        \begin{figure}[H]
            \centering
            \begin{subfigure}{\textwidth}
                \includegraphics[clip, trim=0cm 0cm 0cm 0cm, width=\textwidth]{bsp_gearbox.png}
                \label{fig:bspgearbox}
            \end{subfigure}
        \end{figure}
    \end{minipage}

    \begin{align*}
        J_{\text{red B}} \cdot \omega_B + J_{\text{red C}} \cdot \omega_C &= (J_{\text{red B}} + J_{\text{red C}}) \cdot \omega \\
        n &= \frac{J_{\text{red B}} \cdot n_B + J_{\text{red C}} \cdot n_C}{J_{\text{red B}} + J_{\text{red C}}} = 13{,}31\,s^{-1}
    \end{align*}

    Energieerhaltung mit entstehender Wärme W beim Kuppeln:

    \begin{align*}
        \frac{1}{2} J_{\text{red B}} \cdot \omega_B^2 + \frac{1}{2} J_{\text{red C}} \cdot \omega_C^2 &= \frac{1}{2} (J_{\text{red B}} + J_{\text{red C}}) \cdot \omega^2 + W \\
    \end{align*}

    mit $\omega = 2\pi \cdot n$:
    
    \begin{align*}
        W &= \frac{1}{2} \cdot (2\pi)^2 \cdot \left( J_{\text{red B}} \cdot n_B^2 + J_{\text{red C}} \cdot n_C^2 - (J_{\text{red B}} + J_{\text{red C}}) \cdot n^2 \right) = 3{,}1\,\text{kJ}
    \end{align*}

\end{bspbox}